\paragraph{三角形の二種類の交換}
一つのブロックの長さは \(t\) である。区間の両側の節点の階数を足すと、
中間の二節点の階数の和に等しく、後者の階数差は \(2t\) になる。
負符号の個数も同じ加法関係を満たす。
両側の指標を \(\eta_i\xi^i,\eta_j\xi^j\) とすると、
交換の指標はその商である。選んだ素点で固定指標の \(j-i\) 乗が
非自明であり、捻りの素点でも \(1\le j-i\le m\) の冪が分岐する。
従って交換の \(L_*\) に現れる点は \(1,\ldots,t\) である。
\leanlinks{\leansrc{New/Determinant/Chain/ChainNorms}{階数、符号と指標の確認}}

点データから標準包含のデータに移る計算を記す。有限剰余環 \(R\) 上、
二つの補完行列のブロック幅を \(a',b'\)、可動な上の部分の高さを \(n'\) とする。
重み \(p\) が \(n'\times a'\) 行列の関数であるとき、反対側に渡す重みは
\[
 p_\mathrm{next}\begin{pmatrix}\eta&s\\0&I_{b'}\end{pmatrix}
     =p(-\eta)/|R|^{n'b'}
\]
であり、その他で零である。実際、\(w=\binom{\eta}{0}\) とすれば
\((w\ \ I)\) と \(\binom{I}{-w}\) は同じ可逆行列の逆の行と先頭の列である。
対応する枠の基底をすべて数えると、それぞれ
\(|\mathrm{GL}_{n'+b'}(R)|\)、\(|\mathrm{GL}_{a'}(R)|\) が掛かる。
上側のせん断 \(\begin{psmallmatrix}I&s\\0&I\end{psmallmatrix}\) は行列式が \(1\) なので
核の値を変えない。\(s\) で平均して \(\eta\) を足すとこの重みが得られる。
逆向きの辺では同じ等式を有限 Fourier 変換して使う。
展開すると、せん断の各座標の和は加法指標の直交性により
対応するブロック積の条件を課し、残る正方ブロックでは
原始な行列式指標の行列 Gauss 和になる。
単位元での Gauss 和の絶対値と剰余環の位数が追加される。
これらと \eqref{eq:exchanged-integrals} により、随伴を取る側でも
捻りで変わる絶対値は \(L_*\) と法のノルムの冪である。
\leanlinks{\leansrc{New/Determinant/Chain/SquareLaw}{標準包含のせん断平均と有限 Fourier 変換}
  \leanlinksep \leansrc{New/Determinant/Chain/ChainNorms}{三角形の二つの交換則}}

\paragraph{単独の持ち上げのグラム形式}
階数 \(r\) から \(r+t\) への持ち上げは \(2r<r+t\) を満たす。
この範囲での内積公式は、二つのテータ核の積を
\(W_j\perp W_j^{\mathrm{opp}}\) 上の一つの核にして得られる。
Siegel--Weil の展開をこの核に適用すると、等方な極大部分空間を
走る和が現れる。用いている有限データは、二つの \(r\) 列の
ブロックが指定の原始包含に一致するものに支持される。
従って生き残る等方部分空間は両側への射影が同型で、
有理等長写像のグラフである。グラフ上の零切片は Weil 作用の
行列係数になる。グラフの和を積分の有理点に展開して、
二つの形式の対をその行列係数で右移動した積分に直す。
\leanlinks{\leansrc{Lib/SiegelWeil/RallisInner/Doubling}{倍化した核とグラフの項}}

指定レベルと捻りの素点では原始包含を固定する標準の列条件を置く。
そこでの行列係数の積分は有限データの質量を掛けたレベル平均であり、
固定レベルの形式には恒等写像として作用する。
無限素点の行列係数の積分は、対応する離散系列型の再生性から
零でない固定定数になる。残る有限素点は不分岐であり、
真空の行列係数に重みづけた平行移動の級数 \(R_j\) が残る。
具体的には、レベルとその捻りの法の素点およびそれらの共役の外に
支持された代表元 \(g\) に対し、係数は
\((\chi_{V_j}(\det g)m_f(g)^{r+t})^{-1}\) である。
\(m_f\) は有限真空の Weil 乗数、\(\chi_{V_j}\) は相手空間の指標である。

この級数には単位剰余類の恒等項がある。
不分岐の球で隣接格子への一歩を数え、一歩ごとの個数に
\(q_w^{-(r+t)/2}\) を掛けた重みを \(a_w\) とする。
\(2r<r+t\) のため、有限素点とその共役は、各組で重みの和が
\(1/3\) 以下となるように分けられる。各組の級数で単位項以外の
係数の絶対値和は \(\sum_{k\ge1}(1/3)^k=1/2\) 以下である。
この級数が有界な \(f\) を零に送れば
\(\|f\|_\infty\le\|f\|_\infty/2\) なので \(f=0\) である。
異なる素点に支持された級数は合成で結合し、絶対収束するため、
この単射性が \(R_j\) に移る。持ち上げが零なら、グラム形式の式で
試験形式を \(R_j f\) に取ると \(\|R_j f\|^2=0\) となる。
従って単独の持ち上げは単射である。
\leanlinks{\leansrc{Lib/SiegelWeil/RallisInner/SharpData}{レベル平均と内積公式}
  \leanlinksep \leansrc{Lib/SiegelWeil/UnramifiedSeries/RallisSeries}{真空係数の級数の単射性}}

代数的な形式の内積は、Petersson 内積にレベル指数と Weyl 加群の次元を
掛けたものである。指定レベルでの列条件の質量の分母はこの指数で消える。
捻りの各素点 \(w\) では、階数 \(r\) の部分への射影と残りの \(t\) 列の
自由度から \((\mathrm Nw)^{rt}\) が生じる。正方ブロックの可逆性の条件を
\(|\mathrm{GL}_n(\mathbb F_q)|=q^{n^2}\prod_{s=1}^n(1-q^{-s})\)
で数えると、残る比は \(\prod_{t<s\le t+r}(1-(\mathrm Nw)^{-s})\) になる。
従って捻りで変わるスカラーは次段の \(c_Q(\xi)\) である。
\leanlinks{\leansrc{New/Determinant/Norms/AdjointThetaLift}{有限レベルの質量と Weyl 加群の次元}}

最後に \(W_j\) 上と \(Y_0\) 上の良い形式の比較を行う。
両空間は有限群と無限パラメータを共有し、無限素点の符号が異なる。
良いラベルの組ごとに考えると、\(p_0\) の上の少なくとも一つの分裂素点で
ラベルが原始的である。この成分の中心冪等元は有限一般線型群の
尖点指標から来ており、試験関数のその素点の真放物部分群に沿う定数項は零である。
このため尖点試験関数の跡の比較が使える。
比較で補助素点を使う際には線の安定化群へレベルを下げ、
線に対応する Hecke 作用で剰余体上の特性多項式の零根を単純にする。
その作用の行列式が球ベクトルの個数を与えるので、
元のレベルでの跡の等式に戻る。
従って各球 Hecke 作用素 \(T\) で \(\operatorname{tr}(T^n)\) が等しく、
Newton の恒等式から特性多項式と行列式が一致する。
互いに可換な半単純 Hecke 作用素の有限個で固有系を分離し、
各固有系を抽出する補間多項式の跡を取れば、重複度も一致する。
級数 \(R_j\) の級数和は同じ球 Hecke 作用として両空間に作用するので、
その行列式を \(Y_0\) 上で取れる。
\leanlinks{\leansrc{Lib/Adelic/TraceFormula/CuspidalComparison}{線の安定化群と尖点試験関数の跡の比較}
  \leanlinksep \leansrc{New/Spanning/Multiplicities/EqualMultiplicities}{Hecke 固有系の重複度の一致}}
